% =========================================================
% Archimedean Property
% =========================================================

\subsection{Archimedean Property}

\begin{remark*}[Why this requires proof]
The Archimedean Property states that the natural numbers are unbounded in
$\mathbb{R}$. That conclusion is not a formal consequence of the ordered-field
axioms alone. There exist ordered fields satisfying all ordered-field axioms in
which the natural numbers are bounded above, so the statement must be justified
from stronger structure.

The needed structure is completeness. If $\mathbb{N}$ were bounded above in
$\mathbb{R}$, then the Least Upper Bound Property would produce a supremum
$\alpha \in \mathbb{R}$ for $\mathbb{N}$. But then $\alpha - 1$ cannot be an
upper bound, so there exists $n \in \mathbb{N}$ with $n > \alpha - 1$. Hence
$n+1 > \alpha$, contradicting the fact that $\alpha$ is an upper bound of
$\mathbb{N}$.
\end{remark*}


\begin{theorem}[Archimedean Property]
\label{thm:archimedean-property}
For every real number $x$, there exists a natural number $n$ such that $n > x$.
\end{theorem}

\begin{remark*}[Standard quantified statement]
\[
\forall x \in \mathbb{R}\; \exists n \in \mathbb{N}\; (n > x).
\]
\end{remark*}

\begin{remark*}[Predicate statement]
\[
ArchimedeanProperty.
\]
\end{remark*}

\begin{remark*}[Negated quantified statement]
\[
\exists x \in \mathbb{R}\; \forall n \in \mathbb{N}\; (n \le x).
\]
This means there exists a real number that is an upper bound for all natural numbers. Such a number would place all of $\mathbb{N}$ below a single point of $\mathbb{R}$, which is exactly what the Archimedean Property rules out.
\end{remark*}



\begin{corollary}[Archimedean Corollary]
\label{cor:archimedean-corollary}
For every $x > 0$ and every $y \in \mathbb{R}$ there exists $n \in \mathbb{N}$ such that
$nx > y$.
\end{corollary}

\begin{remark*}[Standard quantified statement]
\[
\forall x \in \mathbb{R}\; \forall y \in \mathbb{R}\;
\bigl(
x > 0 \rightarrow \exists n \in \mathbb{N}\; (nx > y)
\bigr).
\]
\end{remark*}

\begin{remark*}[Negated quantified statement]
\[
\exists x \in \mathbb{R}\; \exists y \in \mathbb{R}\;
\bigl(
x > 0 \land \forall n \in \mathbb{N}\; (nx \le y)
\bigr).
\]
This means there exist a positive real number $x$ and a real number $y$ such that every
natural multiple of $x$ stays below $y$. In other words, the sequence $(nx)$ would be bounded
above despite $x$ being positive.
\end{remark*}

\begin{remark*}[Contrapositive quantified statement]
\[
\forall x \in \mathbb{R}\; \forall y \in \mathbb{R}\;
\bigl(
\forall n \in \mathbb{N}\; (nx \le y)
\rightarrow
x \le 0
\bigr).
\]
This says that if all natural multiples of a real number $x$ remain below a fixed real number
$y$, then $x$ cannot be positive.
\end{remark*}

\begin{lemma}[Integer Part Lemma]
\label{lem:integer-part}
For every real number $x$ there exists a unique integer $m$ such that
\[
m \le x < m+1.
\]
\end{lemma}

\begin{remark*}[Standard quantified statement]
\[
\forall x \in \mathbb{R}\;
\exists m \in \mathbb{Z}\;
\Bigl[
(m \le x < m+1)
\land
\forall k \in \mathbb{Z}\; \bigl((k \le x < k+1) \rightarrow k=m\bigr)
\Bigr].
\]
\end{remark*}


\begin{remark*}[Negated quantified statement]
\[
\exists x \in \mathbb{R}\;
\Bigl[
\forall m \in \mathbb{Z}\; (m > x \lor x \ge m+1)
\;\lor\;
\exists m \in \mathbb{Z}\; \exists k \in \mathbb{Z}\;
\bigl(m \ne k \land m \le x < m+1 \land k \le x < k+1\bigr)
\Bigr].
\]
This means there exists a real number $x$ for which no integer captures $x$ in a half-open
unit interval of the form $[m,m+1)$, or else more than one integer does so. The failure would
lie in the inability to place every real number into a unique adjacent-integer interval.
\end{remark*}




